\documentclass[10pt,a4paper]{amsart}
\usepackage[margin=19mm]{geometry}
\usepackage{amsmath,amssymb,mathtools}
\usepackage[scr=rsfs,cal=euler]{mathalfa}
\usepackage{xfrac}
\usepackage[unicode,hidelinks,bookmarks=false]{hyperref}
\newcommand{\ie}{i.e.\ }
\newcommand{\cf}{cf.\ }
\newcommand{\resp}{respectively\ }
\newcommand{\Subsection}{\subsection}
\providecommand{\nobreakdash}{\nobreak\hbox{-}}
\setlength{\emergencystretch}{2em}
\multlinegap0pt
\usepackage{tikz}
\usetikzlibrary{cd}
\usepackage{mathtools}
\usepackage{float}
\usepackage[shortlabels]{enumitem}
\usepackage{amssymb}
\usepackage{tikz}
\newtheorem{proposition}{Proposition}
\newcommand{\C}{\mathscr{C}}
\newcommand{\NC}{\mathscr{NC}}
\renewcommand{\phi}{\varphi}
\newcommand{\uu}{\mathbb{1}}
\title{Operadic Fibrations and Unary Operadic 2-categories}
\author{Dominik Trnka}
\date{Source: arXiv:2410.05064v2 (CC BY 4.0)}
\begin{document}
\maketitle

\noindent\emph{Source and licence.} Operadic Fibrations and Unary Operadic 2-categories. Version arXiv:2410.05064v2 (CC BY 4.0), distributed by arXiv under the Creative Commons Attribution 4.0 International licence, http://creativecommons.org/licenses/by/4.0/, as recorded in the arXiv licence metadata for this exact version and shown on the abstract page https://arxiv.org/abs/2410.05064v2. Full licence notice: \texttt{licenses/arxiv-2410.05064-CC-BY-4.0.txt}.\par\smallskip

\noindent\textit{Editorial scope.} Counted unit: the article's proposition proposition:unary\_operadic\_2-category\_B (source0.tex lines 782--818). The article's complete original proof of it is delivered with it (source0.tex lines 854--993, one \texttt{proof} environment of 140 lines, which the article places away from the statement). No mathematics is written, repaired or re-derived: the statement and the proof are the article's own lines, in the article's own notation, and no retained line is substituted or re-flowed.  No further source-local context is needed: every cross-reference of every retained line resolves inside the two delivered spans.  Nothing was omitted from any retained span, so the package carries no omission record.  No retained line contains a citation, so the wrapper carries no bibliography and no bibliography entry.  The retained lines contain the article's own diagram code, which the wrapper typesets with the standard tikz packages; no image file is delivered and the package declares no asset dependency.  Editorial formatting only, in addition: an AMS \texttt{amsart} preamble that restates the article's own declarations and macros used by the retained lines, namely the environments \texttt{proposition} on separate counters, together with the standard packages those lines need.  The retained environments print in this document as Proposition 1 in order of appearance, whereas the article prints the same items with the numbering of its own full document.  The complete native source of the article as distributed in the arXiv e-print for this exact version is bundled unchanged as \texttt{sources/166430/source0.tex}.
\begin{proposition}\label{proposition:unary_operadic_2-category_B}
A unary operadic 2-category is a 2-category $\C$ together with the following~data:
\begin{enumerate}
        \item \label{item:1}for every object $a$ of $\C$, there is $\phi(a)\in \pi_0(\C)$,
        \item \label{item:2}for every arrow $x\xrightarrow{q}y$ in $\C$, $\phi(q)$ is an object of $\C$,
        \item \label{item:3}for every lax triangle $g\circ f \xRightarrow{\alpha} h$ in $\C$, $\phi(h)\xrightarrow{\phi(\alpha)}\phi(g)$ is a map in $\C$,
        \item \label{item:4}for a 3-simplex $\sigma_{0123}$ of $\C$,  $\phi(\alpha_{123})\circ \phi(\alpha_{013}) \xRightarrow{\phi(\sigma_{0123})}\phi(\alpha_{023})$ is a lax triangle in $\C$, 
        \item \label{item:5}for every $c \in \pi_0(\C)$, $u_c$ is an object of $\C$,
    \item \label{item:6}for every object $a\in \C$, $a \xrightarrow{u_a} u_{\pi(a)}$ is a map in $\C$,
    \item \label{item:7}for every map $x\xrightarrow{q} y$ of $\C$, there is a 2-cell $u_y\circ q \xRightarrow{u_q} u_x$, and
\item \label{item:8}for every lax triangle $ g\circ f\xRightarrow{\alpha}h$, there is a 3-simplex $u_\alpha=(u_g,u_h,u_f,\alpha)$ in $\C$. 
\end{enumerate}
This data are subject to the following equations. For every $c, a, q, \alpha$, and $\sigma_{0123}$ as above,  
\begin{enumerate} 
\setcounter{enumi}{8}
\item \label{item:9} $ \pi(\phi (q))=\phi(x)$,
    \item\label{item:10} $\pi(u_c)=c$,
    \item\label{item:11}
    $\phi(\uu_a)=u_{\phi(a)}$,
    \item\label{item:12}  $\phi\big( \uu_y\circ q \xRightarrow{\uu_q}q\big)=u_{\phi(q)}$ and $\phi\big(q\circ \uu_x \xRightarrow{\uu_q}q\big)=\uu_{\phi(q)}$,
    
    \item\label{item:13}
   $\phi(\uu,\uu,\alpha,\alpha)=u_{\phi(\alpha)}$,
   $\phi(\uu,\alpha,\alpha,\uu)=\uu_{\phi(g)}\circ\phi(\alpha)\xRightarrow{\uu_{\phi(\alpha)}} \phi(\alpha),$ and
   $\phi(\alpha,\alpha,\uu,\uu)=\phi(\alpha)\circ \uu_{\phi(h)} \xRightarrow{\uu_{\phi(\alpha)}}\phi(\alpha),$
   
    \item\label{item:14} $u_{u_c}=\uu_{u_c}$, $u_{u_a}=\uu_{u_a}$, and $u_{u_q}=\uu_{u_q}$, 
    \item\label{item:15} $\phi\circ u = \uu$, more precisely
    $\phi(u_c)=c$,
    $\phi(u_a)=a$,
     $\phi(u_q)=q$, and
    $\phi(u_\alpha)=\alpha$,
    \item\label{item:16} $\phi(\phi(\sigma_{0123}))=\phi(\alpha_{012})$, $\phi(\phi(\alpha))=\phi(f)$, and $\phi(\phi(q))=\phi(x)$,
    \item\label{item:17} $u_{\uu_x}=\uu_{u_x}$, $u_{s_0 q}=s_0 u_q$, and $u_{s_1 q}=s_1 u_q$. 
    
\end{enumerate}
\end{proposition}

\begin{proof}[Proof of Proposition~\ref{proposition:unary_operadic_2-category_B}]
 It follows from the two lemmas above that a unary operadic 2-category is determined by a diagram 
 % https://q.uiver.app/#q=WzAsNSxbMCwwLCJcXHBpXzBDIl0sWzEsMCwiQ18wIl0sWzIsMCwiQ18xIl0sWzMsMCwiTkNfMiJdLFs0LDAsIk5DXzMiXSxbMSwwLCJcXHBoaSIsMSx7Im9mZnNldCI6M31dLFswLDEsIlxcZXRhIiwxXSxbMSwwLCJcXHBpIiwxLHsib2Zmc2V0IjotM31dLFsyLDEsIlxccGhpIiwxLHsib2Zmc2V0Ijo1fV0sWzEsMiwiXFxldGEiLDEseyJvZmZzZXQiOi0yfV0sWzEsMiwiIiwxLHsib2Zmc2V0IjoyfV0sWzIsMSwiZF8wIiwwLHsib2Zmc2V0IjotNH1dLFszLDIsIlxccGhpIiwxLHsib2Zmc2V0Ijo1fV0sWzIsMV0sWzIsMywiXFxldGEiLDEseyJvZmZzZXQiOi0yfV0sWzMsMl0sWzMsMiwiZF8wIiwwLHsib2Zmc2V0IjotNH1dLFs0LDMsIlxccGhpIiwxLHsib2Zmc2V0Ijo1fV0sWzMsNCwiXFxldGEiLDEseyJvZmZzZXQiOi0yfV0sWzQsM10sWzQsMywiZF8wIiwwLHsib2Zmc2V0IjotNH1dLFsxNSwxNiwiXFxjZG90cyIsMyx7InNob3J0ZW4iOnsic291cmNlIjoyMCwidGFyZ2V0IjoyMH0sImxldmVsIjoxLCJzdHlsZSI6eyJib2R5Ijp7Im5hbWUiOiJub25lIn0sImhlYWQiOnsibmFtZSI6Im5vbmUifX19XSxbMTksMjAsIlxcY2RvdHMiLDMseyJzaG9ydGVuIjp7InNvdXJjZSI6MjAsInRhcmdldCI6MjB9LCJzdHlsZSI6eyJib2R5Ijp7Im5hbWUiOiJub25lIn0sImhlYWQiOnsibmFtZSI6Im5vbmUifX19XV0=
\[\begin{tikzcd}
	{\pi_0\C} & {\C_0} & {\C_1} & {\NC_2} & {\NC_3}
	\arrow["u"{description}, from=1-1, to=1-2]
	\arrow["\phi"{description}, shift right=3, from=1-2, to=1-1]
	\arrow["\pi"{description}, shift left=3, from=1-2, to=1-1]
	\arrow["u"{description}, shift left=2, from=1-2, to=1-3]
	\arrow[shift right=2, from=1-2, to=1-3]
	\arrow["\phi"{description}, shift right=5, from=1-3, to=1-2]
	\arrow["{d_0}", shift left=4, from=1-3, to=1-2]
	\arrow[from=1-3, to=1-2]
	\arrow["u"{description}, shift left=2, from=1-3, to=1-4]
	\arrow["\phi"{description}, shift right=5, from=1-4, to=1-3]
	\arrow[""{name=0, anchor=center, inner sep=0}, from=1-4, to=1-3]
	\arrow[""{name=1, anchor=center, inner sep=0}, "{d_0}", shift left=4, from=1-4, to=1-3]
	\arrow["u"{description}, shift left=2, from=1-4, to=1-5]
	\arrow["\phi"{description}, shift right=5, from=1-5, to=1-4]
	\arrow[""{name=2, anchor=center, inner sep=0}, from=1-5, to=1-4]
	\arrow[""{name=3, anchor=center, inner sep=0}, "{d_0}", shift left=4, from=1-5, to=1-4]
	\arrow["\cdots"{marking, allow upside down}, draw=none, from=0, to=1]
	\arrow["\cdots"{marking, allow upside down}, draw=none, from=2, to=3]
\end{tikzcd}\]
 and simplicial identities up to $X_4=(\NC)_3$.
Most of the identities are satisfied since $\C$ is a 2-category. We translate those which are not automatic:\\
\begin{table}[H]
\resizebox{\textwidth}{!}{%
\begin{tabular}{ccccc}
\begin{tabular}{|c|} 
 \hline
 $\C_1\to \pi_0\C$\\
  \hline
 $\pi\phi=\phi d_0$ \eqref{item:9}\\
 $\phi\phi=\phi d_1$ \eqref{item:16}\\
 \hline
\end{tabular}
&
\begin{tabular}{|c|} 
 \hline
 $\NC_2\to \C_0$\\
  \hline
 $d_0\phi=\phi d_0$ \eqref{item:3}\\
 $d_1\phi=\phi d_1$ \eqref{item:3}\\
 $\phi\phi=\phi d_2$ \eqref{item:16}\\
 \hline
\end{tabular}
&
\begin{tabular}{|c|} 
 \hline
 $\NC_3\to \C_1$\\
 \hline
 $d_0\phi=\phi d_0$ \eqref{item:4}\\
 $d_1\phi=\phi d_1$ \eqref{item:4}\\
 $d_2\phi=\phi d_2$ \eqref{item:4}\\
 $\phi\phi=\phi d_3$ \eqref{item:16}\\
 \hline
\end{tabular}
&
\begin{tabular}{|c|} 
 \hline
 $\C_1\to \NC_3$\\
 \hline
 $u s_0=s_0u$ \eqref{item:17}\\
 $u s_1=s_1u$ \eqref{item:17}\\
 $u u=s_2u$ \eqref{item:14}\\
 \hline
\end{tabular}
&
\begin{tabular}{|c|} 
 \hline
 $\C_0\to \NC_2$\\
 \hline
 $u s_0=s_0u$ \eqref{item:17}\\
 $u u=s_1u$ \eqref{item:14}\\
 \hline
\end{tabular}
\\
\begin{tabular}{|c|} 
 \hline
 $\pi_0\C\to \C_1$\\
 \hline
 $uu=s_0u$ \eqref{item:14}\\
 \hline
\end{tabular}
&
\begin{tabular}{|c|} 
 \hline
 $\pi_0\C\to \pi_0\C$\\
 \hline
 $\pi u=1$ \eqref{item:10}\\
 $\phi u=1$ \eqref{item:15}\\
 \hline
\end{tabular}
&
\begin{tabular}{|c|} 
 \hline
 $\C_0\to \C_0$\\
 \hline
 $\phi s_0=u\phi$ \eqref{item:11}\\
 $d_0u=u \pi$ \eqref{item:6}\\
 $d_1u=1$ \eqref{item:6}\\
 $\phi u=1$ \eqref{item:15}\\
 \hline
\end{tabular}
&
\begin{tabular}{|c|} 
 \hline
 $\C_1\to \C_1$\\
 \hline
 $\phi s_0=s_0\phi$ \eqref{item:12}\\
 $\phi s_1=u\phi$ \eqref{item:12}\\
 $d_0u=u d_0$ \eqref{item:7}\\
 $d_1u=u d_1$ \eqref{item:7}\\
 $d_2u=1$ \eqref{item:7}\\
 $\phi u=1$ \eqref{item:15}\\
 \hline
\end{tabular}
&
\begin{tabular}{|c|} 
 \hline
$\NC_2\to \NC_2$\\
 \hline
 $\phi s_0=s_0\phi$  \eqref{item:13} \\
  $\phi s_1=s_1\phi$  \eqref{item:13}  \\
    $\phi s_2=u \phi$ \eqref{item:13}\\
    $d_0u=u d_0$ \eqref{item:8}\\
 $d_1u=u d_1$ \eqref{item:8}\\
 $d_2u=u d_2$ \eqref{item:8}\\
 $d_3u=1$ \eqref{item:8}\\
 $\phi u=1$ \eqref{item:15}\\
 \hline
\end{tabular}
\end{tabular}%
}
\end{table}

 
 These equations compare directly (cf.~numbers in brackets) to the data and conditions of Proposition~\ref{proposition:unary_operadic_2-category_B}, which finishes the proof.
\end{proof}

\end{document}
